\section{Background and Related Work}
\label{sec:related}

Prior flaky-test repair systems optimize stability without checking fault detection inside the repair loop, whereas test-evolution and empirical-methodology research explains why proxy checks silently erode oracle strength. Table~\ref{tab:related_boundary} positions \textsc{MuGuard-Flake} across three intersecting threads: automated flaky-test stabilization, oracle quality and test flimsiness, and bias-aware, safety-gated testing.

\begin{table*}[t]
\centering
\caption{Methodological positioning of \textsc{MuGuard-Flake} relative to representative flaky test repair, test transformation, and quality-assurance studies.}
\label{tab:related_boundary}
\scriptsize
\setlength{\tabcolsep}{4pt}
\begin{tabular}{>{\raggedright\arraybackslash}p{2.7cm}>{\raggedright\arraybackslash}p{1.8cm}>{\raggedright\arraybackslash}p{4.0cm}>{\raggedright\arraybackslash}p{5.6cm}}
\toprule
Study or System & Venue & Primary Validation Criterion & Limitation or Methodological Role \\
\midrule
\textsc{DexFix}~\cite{zhang2021dexfix} & ICSE 2021 & Template rules and \textsc{NonDex} reruns & Permissive comparators such as non-strict JSON equality mask extraneous or corrupted fields. \\
\textsc{iPFlakies}~\cite{wang2022ipflakies} & ICSE-Comp.\ 2022 & Polluter search and order reruns in Python & Spliced cleaners or minimal setup calls can drop preconditions needed by focal assertions. \\
\textsc{FlakyDoctor}~\cite{chen2024flakydoctor} & ISSTA 2024 & Neurosymbolic localization and LLM reruns & No in-loop fault-detection check verifies whether stabilized patches retain regression sensitivity. \\
\textsc{TestBoost}~\cite{li2024testboost} & ASE 2024 & Fixture-sharing transform and post-hoc mutation audit & Shows 30.7\% of transformed classes lose mutation kills, without an automated repair gate. \\
\textsc{Test Flimsiness}~\cite{parry2026flimsiness} & ICSE 2026 & Multi-run mutation testing for flaky kills & Characterizes noisy kills; we counter flimsiness by requiring consistent kills across repeated isolated runs. \\
Leunig \& Petersen~\cite{leunig2026perl} & SQJ 2026 & Industrial LLM test-repair audit at IBM & Shows green execution conceals assertion dilution; motivates executable conservation gates. \\
Bludau \& Pretschner~\cite{bludau2026szz} & SQJ 2026 & Ground-truth swap for proxy evaluation bias & Motivates measuring the inflation gap between stabilization, in-loop preservation, and external fault retention. \\
Zhang et al.~\cite{zhang2026typerepair} & SQJ 2026 & Present, absent, and mismatched diagnostic control & Inspires our controlled contrast of absent, misleading, and accurate witness feedback. \\
Salva~\cite{salva2026nltestrunner} & SQJ 2026 & False-acceptance vs.\ false-rejection control & Inspires our fail-closed three-state CI gate and explicit pull-request false-rejection audit. \\
\midrule
\textbf{MuGuard-Flake} & \textbf{This Paper} & \textbf{Fail-closed stability and dual-role focal mutation gate, localized witness retry, and witness-backed abstention} & \textbf{Benchmarks against static, dynamic coverage, and cross-vendor LLM-judge gates across five LLMs, historical tools, and developer PRs.} \\
\bottomrule
\end{tabular}
\end{table*}

\subsection{Root Causes and Automated Repair of Flaky Tests}

Intermittent test failures weaken regression signaling in CI by breaking the link between test verdicts and code changes~\cite{luo2014empirical,lam2020lifecycle,parry2021survey}. Empirical taxonomies highlight two context-sensitive categories as primary targets for automated repair~\cite{lam2019idflakies,gyori2016nondex,wang2022ipflakies}. Implementation-Dependent (ID) tests rely on unspecified library behavior, such as hash-map iteration order, reflection ordering, or unordered JSON serialization~\cite{gyori2016nondex,zhang2021dexfix}; \textsc{NonDex} exposes them by permuting internal data structures across seeds. Order-Dependent (OD) tests depend on shared mutable state left behind by earlier tests, such as static fields, singletons, system properties, or files~\cite{lam2019idflakies,shi2019ifixflakies,wang2022ipflakies}. Order-dependent tests split into \emph{victims}, which pass in isolation but fail after a state-polluting test, and \emph{brittles}, which fail in isolation unless preceded by a state-setter test~\cite{shi2019ifixflakies,wang2022ipflakies}.

Automated repair progressed from suite search and syntax templates to LLM agents. \textsc{iFixFlakies}~\cite{shi2019ifixflakies} and \textsc{iPFlakies}~\cite{wang2022ipflakies} mine cleaner or helper statements from the test suite and splice them into order-dependent tests. \textsc{DexFix}~\cite{zhang2021dexfix} applies syntax templates, such as switching unordered sets to linked variants or relaxing exact JSON equality to non-strict mode, to fix implementation-dependent tests. More recently, \textsc{FlakyDoctor}~\cite{chen2024flakydoctor} pairs dynamic localization with iterative LLM repair under rerun validation, and industrial systems such as \textsc{AutoCover}~\cite{rastenis2026autocover} and database flakiness repair pipelines~\cite{berndt2026dbms} adopt similar generate-and-rerun loops.

Across these pipelines, acceptance remains single-sided: any candidate that compiles and passes the stability rerun suite is accepted. Even in template repair, Zhang et al.~\cite{zhang2021dexfix} noted that relaxing JSON checks to ignore key order can admit extra or corrupted fields. Yet neither \textsc{FlakyDoctor}~\cite{chen2024flakydoctor} nor suite-search tools verify inside the repair loop whether a stabilized test preserves its original fault-detection power.

\subsection{Patch Overfitting, Oracle Quality, and Test Flimsiness}

Oracle erosion in flaky test repair is a sharper form of patch overfitting in automated program repair (APR). In production APR, generate-and-validate loops often emit plausible patches that pass weak test suites by deleting behavior or bypassing faulty paths~\cite{qi2015kali,smith2015cure}. Test repair faces a deeper structural asymmetry~\cite{barr2015oracle,daniel2009reassert,li2019trip,rahman2025utfix}. Whereas production APR leaves the test suite intact as an external referee, flaky test repair edits the test oracle and fixture themselves. Once an assertion is weakened, CI reports a permanently green build with no failure signal.

Prior test-evolution studies confirm that automated test edits can silently mask regressions, yet existing checks leave an online verification gap. Early tools such as \textsc{ReAssert}~\cite{daniel2009reassert} risked hiding bugs by rewriting expected values from live runs. More recently, \textsc{TestBoost}~\cite{li2024testboost} found post hoc that 30.7\% of fixture-transformed classes lost mutation kills, and change-aware repair tools~\cite{li2019trip,rahman2025utfix} target evolving production code rather than flaky-test stabilization. Dynamic coverage and static diff review also fall short as online guards. Checked Coverage~\cite{schuler2013checked} demonstrated that statement and branch coverage record line execution rather than whether assertions inspect computed state, and standalone LLM judges score diff plausibility without observing which production branches stop influencing test verdicts.

Evaluating test edits with mutation testing builds on extensive evidence linking mutant kills to real fault detection~\cite{just2014mutants,papadakis2019mutation}, yet introduces its own measurement hazards. Automatically generated suites are frequently flaky~\cite{gruber2024generated}, and mutation scores in LLM-generated suites shift with granularity and code version~\cite{zhao2026correlation}. Most directly, Parry et al.~\cite{parry2026flimsiness} formalized \emph{Test Flimsiness}: mutating production code can trigger intermittent failures in otherwise passing tests, causing single-run mutation testing to credit chance failures as genuine kills. To rule out flimsy kills, Section~\ref{sec:construct} retains only baseline mutants that fail deterministically in the target test across repeated isolated runs, and cross-checks accepted repairs on disjoint held-out faults.

\subsection{Proxy Bias, Diagnostic Controls, and Safety-Gated Testing}

Sound evaluation of repair gates must confront two recurring validity threats: green pass rates that conceal weakened assertions and proxy metrics that conflate distinct quality dimensions. In an industrial study at IBM, Leunig and Petersen~\cite{leunig2026perl} showed that pass rates hide assertion dilution during iterative LLM test repair, and Bludau and Pretschner~\cite{bludau2026szz} demonstrated how noisy proxy labels distort empirical rankings. Following those lessons, we separate stability, in-loop fault preservation, independent held-out fault detection, and semantic erosion types instead of collapsing them into a single repair rate.

Turning diagnosis into a practical CI gate also requires testing whether error feedback truly drives repair gains and measuring the cost of conservatism. Adopting the diagnostic-control logic of Zhang et al.~\cite{zhang2026typerepair}, who contrasted accurate, absent, and mismatched error signals in LLM repair, we test whether localized surviving-mutant feedback outperforms bare numeric counts or misleading hints. On the deployment side, Salva~\cite{salva2026nltestrunner} framed LLM testing agents around explicit retry-and-abstention control of false acceptances versus false rejections. \textsc{MuGuard-Flake} unifies these principles in a fail-closed CI merge gate whose false-rejection rate, external fault-detection benefit, and runtime overhead we measure end to end. Closing the gaps across these three threads first requires formalizing oracle erosion and non-circular external validation, which we turn to next.
